\phantomsection
\addcontentsline{toc}{section}{Resumen}
\begin{center}\bfseries Resumen\end{center}
\begingroup
\fontsize{9.5}{10.7}\selectfont
\setlength{\parskip}{1.5pt plus .5pt minus .5pt}
\noindent
Se derivan las distribuciones de Bose--Einstein y Fermi--Dirac, su límite
diluido de Maxwell--Boltzmann y las leyes de radiación de Planck,
Stefan--Boltzmann y Wien a partir de las realizaciones cuánticas y térmicas
construidas en Holografía Modular Triádica y Mecánica Dimensional.
El soporte aritmético de caminos, denominado \emph{All Possible Paths}
(APP), es el toro discreto \((\mathbb Z/9\mathbb Z)^2\), con grafo de
Cayley aditivo de direcciones cardinales y dos evaluaciones, suma y
producto. La firma ternaria orientada, TRIT, toma valores en
\(\{-1,0,+1\}\) y distingue orientación y régimen; el núcleo topológico
de fase, \emph{Topological Phase Kernel} (TPK), compone selección,
transporte, actualización y prolongación conservando la memoria.
Esta dinámica produce los estados y registros regionales sobre los que
se definen la evolución unitaria y la acción del calendario.

La paridad de las realizaciones determina sus completaciones simétrica y
exterior: la primera admite ocupaciones enteras no negativas; la segunda
produce la exclusión algebraica de Pauli. Las multiplicidades determinan los recuentos microcanónicos; la
factorización del estado de Gibbs da las dos distribuciones de equilibrio,
cuyo límite diluido se cuantifica mediante cotas explícitas. Para los
modos electromagnéticos transversales libres, con potencial químico nulo,
la ocupación bosónica y la densidad de modos producen el espectro de
Planck. Se demuestra la convergencia de las sumas discretas mediante
cotas uniformes en bajas y altas frecuencias. Sus momentos determinan
las densidades de número, energía y entropía; la integración angular da
el flujo de Stefan--Boltzmann, y la extremización de las densidades en
frecuencia y longitud de onda da las dos leyes de Wien.

El calendario y la holonomía producen además una renovación de capacidad
$(20,24,25,23)$, un portador graduado de multiplicidades
$1/380/649$ y un espectro geométrico nativo de razón $1/10$. Su
refinamiento trítico conserva el residuo y da $1/1000$; la sección marcada
$23/26/29$ evalúa exactamente
$1{,}380649\times10^{-23}$. El origen, el marcador y la procedencia
energética se transportan mediante isometrías explícitas. Las corrientes de
borde fijan el desfase sectorial, y la composición con la acción demuestra
$G_a k_B\Theta_{{\rm clk},a}\log3=c^4L_\alpha$ y la identidad efectiva del
transductor sobre su sección positiva.

La entropía de configuraciones, la medida de historias y el espectro de
un estado reducido quedan relacionados por mapas de transporte y
reducción. Una identidad de divergencia relativa caracteriza la medida
estacionaria de entropía máxima del envolvente areal--volumétrico.
La transducción térmica enlaza información, acción y calendario con sus
bases dimensionales. Las condiciones de representación, equilibrio y
convergencia especifican el dominio de cada resultado.

En la realización finita \(\Lambda=(\mathbb Z/27\mathbb Z)^2\), la
transformada de Fourier unitaria establece
\(|\operatorname{supp}\psi|\allowbreak\,|\operatorname{supp}\widehat\psi|\ge729\)
y \(H_{\rm pos}(\psi)\allowbreak+H_F(\psi)\ge\log729\); cuatro colecciones subgrupales
saturan ambas cotas. En el dominio dodecafásico, los lectores de F044
distinguen correlaciones de orden que la entropía marginal no conserva.
El observador primo de vacancias \(R_{\rm vac}\) compone la incidencia
esturmiana con el reloj logarítmico y satisface una identidad modal por
sumas simétricas o de Fejér. Sus seis cortes hasta \(10^8\) y los treinta
modos registrados son testigos finitos; la reproducción independiente
alcanza \(10^6\). La estimación asintótica crítica requiere cotas modales
uniformes en \(k\).

\medskip
\noindent\textbf{Palabras clave:} Bose--Einstein; Fermi--Dirac; ley de Planck;
estados genealógicos; completación de Fock; dinámica de Floquet;
entropía de ocupación; Hamiltoniano modular; límite termodinámico;
radiación térmica; incertidumbre entrópica; vacancias primas; memoria
ordenada.
\par\endgroup
